
\documentclass[12pt, a4paper, twoside]{article}

%--Packages--
\usepackage[english]{babel}
\usepackage[utf8]{inputenc}
\usepackage[T1]{fontenc}
\usepackage{showkeys}

\usepackage{latexsym,amsmath,amsthm,amssymb,bm,lmodern,titlesec,xcolor,graphicx,setspace,pdfpages,
geometry,marginnote,caption,array,sectsty,siunitx,eurosym,colortbl,float,wrapfig,vmargin,tikz,hyperref}
\usepackage[nottoc]{tocbibind}
\usepackage{tikz-cd}


\setpapersize{A4}
\setmarginsrb{30mm}{20mm}{30mm}{30mm}%left top right bottom
                 {8mm}{6mm}{10mm}{10mm}%headheight headsep footheight footskip

% Some useful colors
\definecolor{Rosso}{RGB}{232,61,61}
\definecolor{Viola}{RGB}{188,30,188}
\definecolor{Celeste}{RGB}{34,186,211}
\definecolor{Arancione}{RGB}{232,181,19}
\definecolor{Blu}{RGB}{0, 44, 92}
\definecolor{Azzurro}{RGB}{57,186,238}


\hypersetup{colorlinks=true,linkcolor=Blu,allbordercolors=white, urlcolor=blue, citecolor=Viola}
\usepackage{enumitem} 
\renewcommand\labelitemi{\textbullet} 


% math environments
%  we use \newtheorem to specify our theorem-like environments
   % (theorem, definition, etc.) and how to display and number them.
   %
   % Note: The \theoremstyle declarations affect the appearance of the
   % Theorems, Definitions, etc.
   % Note 2: If you want swedish names you have to go and change the argument in brackets e.g. \newtheorem{thm}[defi]{Theorem} should become \newtheorem{thm}[defi]{Sats}
\theoremstyle{definition}
\newtheorem{definition}{Definition}[section]

\newtheorem{hottake}{Hot Take}[section]

\newtheorem*{problem}{Problem}
\newtheorem{conjecture}[definition]{Conjecture}

%\theoremstyle{plain}
\newtheorem{theorem}[definition]{Theorem}
\newtheorem{proposition}[definition]{Proposition}
\newtheorem{corollary}[definition]{Corollary}
\newtheorem{lemma}[definition]{Lemma}
%\newtheorem*{remark}{Remark}%[definition]{Remark}
%\newtheorem{example}[definition]{Example}
%\newtheorem*{notation}{Notation}
%\newtheorem*{note}{Note}

\theoremstyle{remark}
\newtheorem*{question}{Question}
\newtheorem{remark}[definition]{Remark}
\newtheorem{example}[definition]{Example}
\newtheorem*{notation}{Notation}
\newtheorem*{note}{Note}

%\onehalfspacing
\numberwithin{equation}{section}

%--Heading of  section--
\titleformat{\section}[display]
  {\normalfont\bfseries\color{Blu}}
  {\filleft
    \begin{tikzpicture}
    \node[
      outer sep=0pt,
      text width=2.5cm,
      minimum height=3cm,
      fill=white,
      font=\color{Blu}\fontsize{80}{0}\selectfont,
      align=center
      ] (num) {\thesection};
    \node[
      rotate=90,
      anchor=south,
      font=\color{black}\large\normalfont
      ] at ([xshift=-5pt]num.west) {\Huge\textsc{}};  
    \end{tikzpicture}%
  }
  {10pt}
  {\titlerule\vskip4pt\Huge\sc\sffamily}

%--Heading of subsection--
\newcommand*\numb[1]{%
\begin{tikzpicture}[baseline=-0.7ex]
\node[
  outer sep=0pt,
      text width=0.6cm,
      minimum height=0.6cm,
      fill=Blu,
      font=\color{white}\fontsize{12}{20}\selectfont,
      align=center
      ] (num) {\thesubsection};
\end{tikzpicture}%
}
\titleformat{\subsection}
  {\normalfont\color{Blu}\large\sc\sffamily}{\numb{\thesubsection}}{0.8em}{}

%--Heading subsubsection--
\newcommand*\stocaz[1]{%
\begin{tikzpicture}[baseline=-0.7ex]
\node[
  outer sep=0pt,
      text width=0.8cm,
      minimum height=0.5cm,
      fill=Blu,
      font=\color{white}\fontsize{10}{20}\selectfont,
      align=center
      ] (num) {\thesubsubsection};
\end{tikzpicture}%
}
\titleformat{\subsubsection}
  {\normalfont\color{Blu}\bfseries\sc\sffamily}{\stocaz{\thesubsubsection}}{0.8em}{}
  
% Blackboard Boldface
\newcommand{\NN}{\mathbb{N}}
\newcommand{\ZZ}{\mathbb{Z}}
\newcommand{\QQ}{\mathbb{Q}}
\newcommand{\RR}{\mathbb{R}}
\newcommand{\CC}{\mathbb{C}}
\newcommand{\KK}{\mathbb{K}}
\newcommand{\FF}{\mathbb{F}}

\newcommand{\1}{\mathbf{1}}
\newcommand{\0}{\mathbf{0}}

% Text
\newcommand{\GL}{\text{GL}}
\newcommand{\SL}{\text{SL}}
\newcommand{\SymP}{\text{Sym}}
\newcommand{\AltP}{{\bigwedge}}
\newcommand{\Hom}{\text{Hom}}
\newcommand{\SO}{\text{SO}}
\newcommand{\SU}{\text{SU}}
\newcommand{\id}{\text{id}}
\newcommand{\Stab}{\text{Stab}}
\newcommand{\Cent}{\text{Cent}}
\newcommand{\Norm}{\text{Norm}}
\newcommand{\Type}{\text{Type}}
\newcommand{\Dih}{\text{D}}
\newcommand{\Sym}{\mathcal{S}} %{\mathfrak{S}} % \text{S}
\newcommand{\Alt}{\text{A}}
\newcommand{\Klein}{\text{V}}
\newcommand{\Cyc}{\mathcal{C}}%{\text{C}}
\newcommand{\sgn}{\text{sgn}}
\newcommand{\Span}{\text{Span}}
\newcommand{\im}{\text{im }}
\newcommand{\Tr}{\text{Tr }}
% Vectors
\newcommand{\bas}{\mathbf{\hat{e}}}

\newcommand{\ihat}{\mathbf{\hat{\imath}}}
\newcommand{\jhat}{\mathbf{\hat{\jmath}}}
\newcommand{\ghat}{\mathbf{\hat{g}}}
\newcommand{\hhat}{\mathbf{\hat{h}}}
\newcommand{\vhat}{\mathbf{\hat{v}}}
\newcommand{\what}{\mathbf{\hat{w}}}
\newcommand{\uhat}{\mathbf{\hat{u}}}
\newcommand{\xhat}{\mathbf{\hat{x}}}

\newcommand{\vvec}{\mathbf{v}}
\newcommand{\uvec}{\mathbf{u}}
\newcommand{\wvec}{\mathbf{w}}


\newcommand{\Ss}{\scriptstyle}

\newcommand{\barpi}{\bar{\pi}}



\begin{document}

% This includes the front page (you need to compile the cover.tex separately and have the .pdf available in the same folder as this .tex file).
%\setlength{\voffset}{0cm}\setlength{\hoffset}{0cm}\includepdf[pages={1,2,3,4}]{cover}\setlength{\voffset}{-2.54cm}\setlength{\hoffset}{-2.54cm}



%\clearpage{\thispagestyle{empty}\cleardoublepage}

%\begin{abstract}
%    Your summary goes here
%\end{abstract}

%create the table of contents
\clearpage{\thispagestyle{empty}\cleardoublepage}
{\hypersetup{linkcolor=black}
\tableofcontents
}

%% your first section
\clearpage{\thispagestyle{empty}\cleardoublepage}
\section{Introduction}

Lorem ipsum dolor sit amet, consectetur adipiscing elit, sed do eiusmod tempor incididunt ut labore et dolore magna aliqua.
Ut enim ad minim veniam, quis nostrud exercitation ullamco laboris nisi ut aliquip ex ea commodo consequat. 
Duis aute irure dolor in reprehenderit in voluptate velit esse cillum dolore eu fugiat nulla pariatur. 
Excepteur sint occaecat cupidatat non proident, sunt in culpa qui officia deserunt mollit anim id est laborum.

You can also use swedish letters ä ö å, Ä, Ö, Å. Unicode characters in general should work.


%Another section
\clearpage{\thispagestyle{empty}\cleardoublepage}
\section{Theorems and definitions}

\begin{definition}[Test av definition]
	Test av definition
\end{definition}

\begin{problem}[Test av problem]
	Test av problem
\end{problem}

\begin{conjecture}[Test av conjecture]
	Test av conjecture
\end{conjecture}

\begin{theorem}[Test av theorem]
	Test av theorem
\end{theorem}

\begin{proposition}[Test av proposition]
	Test av proposition
\end{proposition}

\begin{corollary}[Test av corollary]
	Test av corollary
\end{corollary}

\begin{lemma}[Test av lemma]
	Test av lemma
\end{lemma}

\begin{question}[Test av question]
	Test av question
\end{question}

\begin{remark}[Test av remark]
	Test av remark
\end{remark}

\begin{example}[Test av example]
	Test av example
\end{example}

\begin{notation}[Test av notation]
	Test av notation
\end{notation}

\begin{note}[Notering]
	This note's for you
\end{note}


\begin{definition}
	This is a definition!
\end{definition}

\begin{theorem}\label{thm: label}
	This is a theorem!
\end{theorem}

To cross reference: By Theorem~\ref{thm: label} we can say...
\begin{lemma}
	This is a lemma.
\end{lemma}
\begin{proof}
	Put here the proof of the lemma!
\end{proof}


\begin{proof}[Proof of Theorem \ref{thm: label}]
	Another proof
\end{proof}


\begin{corollary}
	here is a corollary!
\end{corollary}

If you want to write a beautiful equation in line do like that $e^{i\pi}+1=0$.
You might also want to write it on its own line
\[
e^{i\pi}+1=0.
\]
If you think that you are going to need to cross reference your equation:
\begin{equation}\label{eq:beautiful}
	e^{i\pi}+1=0
\end{equation}
and this is how you cross reference it \eqref{eq:beautiful}
You might also want to have aligned equations on more lines.
\begin{align}
	e^{i\pi}+1&=0\\
	(\cos\theta)^2+(\sin\theta)^2 &=1
\end{align}

If you need to cross reference something in the bibliography, use the cite command: \cite[p. 45]{Biggs2002}. 
The actual references are stored in a different file, called bibliographyFile.bib 

You can reference most things you found useful,
for example, books \cite{Biggs2002,Knuth1998ArtOfProgramming,LascouxPolynomialsBook}.
Research articles and preprints, \cite{AlexanderssonSulzgruber2019,Sherman2015x}.
Or another thesis, \cite{Teff2013}.

You can also cite online sources. For example, the Encyclopedia of integer sequences, \cite{OEIS},
or a YouTube video \cite{Bazett2018yt}. 
It is even possible to refer to a discussion you had with someone, \cite{PrivateCommunicationWithJustin}.
This is usually rare, it is better to thank them in the acknowledgements section (if they were helpful).



%another section
\clearpage{\thispagestyle{empty}\cleardoublepage}
\section{Tables and graphs}
Here a few examples of tables and graphs.
%A subsection
\subsection{Tables}
\begin{center}
	\begin{tabular}{c c c c c c c c}
		\arrayrulecolor{Azzurro}
		\hline
		{\bfseries Code} & {\bfseries CdL} & {\bfseries xxx} & {\bfseries $T_{setup/lotto}$} & {\bfseries $T_{lav/st}$} & {\bfseries $T_{proc/pezzo}$} & {\bfseries Quantitity} & {\bfseries $T_{tot}$}\\
		\hline
		100 & 4 & 250 & 25 & 0,5 & 0,6 & 1 & 0,6\\
		111 & 2 & 250 & 20 & 2 & 2,08 & 1 & 2,08 \\
		111 & 3 & 250 & 15 & 1,5 & 1,56 & 1 & 1,56 \\
		112 & 2 & 250 & 20 & 2,5 & 2,58 & 1 & 2,58 \\
		112 & 3 & 250 & 15 & 2 & 2,06 & 1 & 2,06\\
		113 & 3 & 500 & 15 & 1 & 1,03 & 2 & 2,06\\
		120 & 1 & 50 & 30 & 2 & 2,6 & 0,1 & 0,26\\
		121 & 1 & 25 & 30 & 3 & 4,2 & 0,1 & 0,42 \\
		121 & 1 & 25 & 30 & 2,5 & 3,7 & 0,1 & 0,37 \\
		\hline
	\end{tabular}
\end{center}



%another section
\clearpage{\thispagestyle{empty}\cleardoublepage}

\section{Insert a figure}
\begin{figure}[H]
	\centering
	\includegraphics[width=1\textwidth]{grafo.png}
	\caption{\label{fig:1} put the caption here}
\end{figure}

\subsection{Footnote}

You can create a footnote like this.\footnote{I created a footnote.}



\clearpage{\thispagestyle{empty}\cleardoublepage}
\section{Conclusion}

Remember the internet is full of places that can help you write a beautiful document with latex.
Some example are \url{overleaf.com} and \verb!stackexchange! just ask Google! 
In addition with an SU e-mail address you should be able to 
create (for free) a premium account on \url{overleaf.com}, 
so you do not need to download a tex editor and you will have
the full latex distribution at your hand (if you are connected to the internet).


$$ \NN \ \ZZ \ \QQ \ \RR \ \CC \ \KK \ \FF$$
$$ \1 \ \0 $$

$$ \GL(V),  \SL_n(\KK), \SymP^k V, \AltP^k V, \Hom(V,W), \SO, \SU, \id, \Stab(g), \Cent(g), \Norm(g),$$ $$ \Type(\sigma), \Dih_{2n}, \Sym_n, \Alt_n, \Klein, \Cyc_n, \sgn(\sigma), \Span(\bas_i)_i $$

$$ \bas \ihat \jhat \ghat \hhat \vhat \what \uhat \vvec \uvec \wvec $$
\clearpage{\thispagestyle{empty}}
\section{Preliminary topics}

\subsection{Linear algebra topics}

	\paragraph{Definition of kernel and image of a map}
	
	\begin{definition}[Kernel and image of a linear map]\label{def:kernelimage}
		Let $V$ and $W$ be two vector spaces and let $\varphi: V \rightarrow W$ be a linear map. Then the kernel and the image of the map are defined thusly:
		\begin{itemize}
			\item[i)] $\ker \varphi = \left\lbrace \vvec \in V \ \middle\vert \ \varphi(\vvec) = \0 \right\rbrace = \varphi^{-1}(\0)$.
			\item[ii)] $\im \varphi = \left\lbrace \wvec \in W \ \middle\vert \ \exists \vvec \in V \text{ s.t. } \varphi(\vvec) = \wvec \right\rbrace = \varphi(V)$.
		\end{itemize}
	\end{definition}
	
	\begin{remark}
		The kernel and the image of a linear map are subspaces of the domain and codomain of the map respectively, ie. $\ker \varphi$ is a subspace of $V$ and $\im \varphi$ is a subspace of $W$.
	\end{remark}
	
	\paragraph{Existence of complementary subspaces}
	
	\begin{corollary}\label{thm:compsubspaces}
		Let $V$ be a vector space and $W$ be a vector subspace of $V$. Then there exists a complementary vector subspace $W'$ in $V$ such that $W \cap W' = \emptyset$ and $W \cup W' = V$. This is equivalent to saying that $V$ is the direct sum of $W$ and $W'$, denoted as $V = W \oplus W'$.
	\end{corollary}
	\begin{proof}
		La la la
	\end{proof}
	
\subsection{Group theoretical topics}

	\subsubsection{Group action}
	
		We could let the group $G$ act on a set $X$ and study the action of $G$ on $X$ defined by
		\begin{align*}
			G \times X \rightarrow X
		\end{align*}
		and thereby permuting the elements of $X$ by for any $g \in G$,
		\begin{align*}
			(g,x) \mapsto g \cdot x = gx,
		\end{align*}
		for any $x \in X$.
	
	\subsubsection{The symmetric group}
	
		Recall that two elements $\sigma, \tau \in \Sym_n$ are conjugate if there exists a $\pi \in \Sym_n$ such that $\sigma = \pi \tau \pi^{-1}$. The conjugacy classes partition $\Sym_n$ into disjoint subsets. 
		
		\begin{proposition}[Conjugacy in $\Sym_n$]
			Any elements $\sigma, \tau \in \Sym_n$ are conjugate if and only if they are of the same cycle type.
		\end{proposition}
		
		\begin{definition}[Sign of a permutation]
			\begin{align*}
				\sgn(\sigma) = (-1)^k,
			\end{align*}
			where $k$ is the number of transpositions required to compose $\sigma$.
		\end{definition}
		
		\begin{proposition}\label{prop:signwelldefined}
			The sign of a permutation is well-defined. That is a permutation is either \textit{even} with sign +1 or \textit{odd} with sign -1.
		\end{proposition}
		\begin{proof}
			\cite[Thm.12.6.1.]{Biggs}
		\end{proof}
		
		
		

\clearpage{\thispagestyle{empty}}
\section{Representations of finite groups}

\marginnote{Include note on source!} 

\paragraph{First definitions.}

The main purpose of this text is to study linear representations of finite groups.

Let $G$ be a finite group written multiplicatively. Denote its order by $|G|$ and the action of its elements with $\cdot$, eg. $g \cdot x$.

Let $V$ be a finite-dimensional vector space over the field of complex numbers $\CC$. Its group of automorphisms is denoted by $\GL(V)$. If $V$ is provided with a basis, usually denoted $(\bas_i)_{i=1}^n$, where $n = \dim V$, then $\GL(V)$ is identified with the set of linear transformations of $V$, denoted $\GL_n(\CC)$\cite[18.1]{DummitFoote}.

\begin{definition}[Representation]\cite{Serre}
	Given a finite group $G$, a representation of $G$ in the vector space $V$ is a homomorphism
	\begin{align}\label{eq:DefRepr}
		\rho: G \rightarrow \GL(V).
	\end{align}
\end{definition}
For every $g$ in $G$, there is a linear map $\rho_g$ in $\GL(V)$ with an action on the elements of $V$ defined by 
\begin{align*}
	g \cdot \vvec = \rho_g (\vvec).
\end{align*}
The linearity of the map means that for any $\vvec, \wvec \in V$,
\begin{align}\label{eq:linearityOfg}
	g \cdot (a\vvec + b\wvec) = ag\cdot\vvec + bg\cdot\wvec,
\end{align} 
and also since it is a homomorphism, for any $g,h \in G$ we have
\begin{align}\label{eq:homomorphism}
	\rho_{gh} = \rho_g \rho_h.
\end{align}
From the homomorphism of $\rho$, two consequences follow.
\begin{proposition}
	For the identity element $e$ of $G$ and an arbitrary element $g$ with inverse $g^{-1}$ in $G$ we have,
	\begin{align}\label{eq:rhoIdAndInv}
		\rho_e = \id \quad \text{and} \quad \rho_g^{-1} = \rho_{g^{-1}},
	\end{align}
	where $\id$ is the identity transformation and $\rho_{g^{-1}}$ is the inverse of the representation of $g$.
\end{proposition}
\begin{proof}
	Take $g$ as an arbitrary element of $G$. The first identity follows from taking $h=e$ in Equation~\ref{eq:homomorphism}:
	\begin{align*}
		\rho_{eg} = \rho_e \rho_g.
		%\rho(eg) = \rho(e)\rho(g). %	\rho(g) \overset{!}{=} \rho(eg) = \rho(e)\rho(g)
	\end{align*}
	Since $eg=g$ for any $g\in G$ we must have that $\rho_g = \rho_e\rho_g$ which is true if and only if $\rho_e = \id$. Instead taking $h=g^{-1}$ we have:
	\begin{align*}
		\rho_{gg^{-1}} = \rho_g \rho_{g^{-1}},
		%\rho(g) \rho(g^{-1}) =  \rho(gg^{-1}~). %\overset{!}{=} \rho(e) = \id.
	\end{align*}
	but $gg^{-1} = e$ and $\rho_e = \id$, implying $\rho_{g^{-1}} = \rho_g^{-1}$.
\end{proof}

\begin{note}
	A vector space $V$ together with a group action fulfilling Equations \ref{eq:DefRepr}, \ref{eq:homomorphism} and \ref{eq:rhoIdAndInv} gives $V$ the structure of a $G$-module \cite[1.3]{Sagan}. To avoid a treatment on modules, if a vector space $V$ is a $G$-module it will be referred to as a \textit{representation space} of $G$.
\end{note}

\begin{note}
	The homomorphism $\rho$, the set of maps $\{\rho_g\}_{g \in G}$ and the representation space $V$ are interchangeably and abusively called the \textit{representation of $G$}.
\end{note}

\paragraph{Choice of basis.} 

If a basis $\{\bas_i\}_{i=1}^n$ is provided for a representation space, then a \textit{matrix representation} can be provided. In this case, $\rho_g$ is the atrix in $\GL_n(\CC)$ associated with the action of $g$ expressed in the basis provided. A matrix representation is not canonical and is dependent on the basis chosen.

\paragraph{Degree of a representation.} 

\begin{definition}[Degree of a representation]
	Let $V$ be a representation space of $G$. Then the dimension of $V$ is referred to as the degree of the representation.
\end{definition}	
\subsection{Examples}\label{sect:basicreprs}

This section follows \cite[1.2]{Serre}.

\subsubsection{Trivial representations}

%	A representation $\rho$ where every $g$ is such that $|\rho_g| = 1$, is called a unital representation.

\begin{example}[Trivial representation]\label{example:trivrepr}
	For any group $G$ there is a trivial degree 1 representation defined by mapping each element of $G$ to 1, ie.
	\begin{align*}
		\rho_g = 1,
	\end{align*}
	for every $g \in G$. It is clearly a homomorphism since for any $g,h \in G$ we have that
	\begin{align*}
		\rho_{gh} = 1 = 1 \cdot 1 = \rho_g \rho_h.
	\end{align*}
	This could be extended to any vector space by mapping $g$ to the identity transformation of that vector space\footnote{For a vector space of dimension $n$, the representation taking every $g$ in a group to the $n \times n$ identity matrix can be described, in the language of Section~\ref{sect:tensorrepr}, as the direct sum of $n$ copies of the trivial representation.}.
\end{example}

\begin{note}
	The trivial representation (the mapping is trivial) is not to be confused with the trivial space (the vector space containing only the zero vector).
\end{note}

\begin{example}[Alternating representation of $\Sym_n$]\label{example:altrepr}
	Choosing $G = \Sym_n$, another degree 1 representation can be found by studying the signs, or parities, of the elements of $\Sym_n$. By \cite[Thm.12.6.1.]{Biggs}, the sign of a permutation is well-defined, so for any two permutations $\sigma$ and $\tau$ with respective signs $(-1)^s$ and $(-1)^t$, their composition has the sign
	\begin{align*}
		\sgn(\sigma\tau) = (-1)^{s+t} = (-1)^s \cdot (-1)^t = \sgn (\sigma) \cdot \sgn (\tau),
	\end{align*}
	so clearly the map $\sgn: \Sym_n \rightarrow \{\pm 1\}$ is a homomorphism and thus a representation of degree 1, where even permutations are mapped to $+1$ and odd to $-1$.
\end{example}

\subsubsection{Degree 1 representations of $\Cyc_n$}\label{sect:deg1cycn}

Choose $G = \Cyc_n$, and let $g$ be a generator of $\Cyc_n$ such that 
\begin{align*}
	\Cyc_n = \{ e, g, g^2, \dots, g^{n-1}\}
\end{align*}
and $g^n = e$. Consider a map $\rho: \Cyc_n \rightarrow \CC$ defined as a homomorphism by $\rho_{g^a}\rho_{g^b} = \rho_{g^{a+b}}$ for some integers $a,b$, then by Equation~\ref{eq:rhoIdAndInv} we have that $\rho_e = 1$, but $g^n = e$ and by induction on Equation~\ref{eq:homomorphism} we have that $\rho_{g^n} = (\rho_g)^n$, so then we must have that $(\rho_g)^n = 1$, that is $\rho_g$ is mapped to a $n$th root of unity.
In conclusion, for $\Cyc_n$ we have found $n$ representations of degree 1, each mapping $g$ to a $n$th root of unity and the powers of $g$ to the corresponding powers of the root of unity. 

\begin{example}[$\Cyc_3$]
	The three third roots of unity are $1$, $\omega = \frac{-1+i\sqrt{3}}{2}$ and $w^2 = \frac{-1-i\sqrt{3}}{2}$, so the three representations of $\Cyc_3$ are presented in table \ref{table:Cyc3}. Note that $\rho_1^* = \rho_2$.
	
%	\begin{table}[hbt!]
%		\begin{tabular}{c | c c c}
%			$\Cyc_3$ & $e$ & $g$        & $g^2$      \\ \hline
%			$\rho_0$          & 1   & 1          & 1          \\
%			$\rho_1$          & 1   & $\omega$   & $\omega^2$ \\
%			$\rho_2$          & 1   & $\omega^2$ & $\omega$
%		\end{tabular}
%		\centering
%		\caption{Three representations of $\Cyc_3$. $\omega = e^{2 \pi i/3}$}
%		\label{table:Cyc3}
%	\end{table}
\end{example}

\begin{example}[$\Cyc_4$]
	The four fourth roots of unity are $1,i,-1$ and $-i$ and the four corresponding representations of $\Cyc_4$ are presented in Table~\ref{tbl:cyc4}. Note that $\rho_1^* = \rho_3$.
%	\begin{table}[hbt!]
%		\centering
%		\begin{tabular}{c | c c cl}
%			$\Cyc_4$ & $e$ & $g$  & $g^2$ & $g^3$ \\ \hline
%			$\rho_0$           & 1   & 1    & 1     & 1     \\
%			$\rho_1$           & 1   & $i$  & $-1$  & $-i$  \\
%			$\rho_2$           & 1   & $-1$ & $1$  & $-1$   \\
%			$\rho_3$           & 1   & $-i$ & $-1$   & $i$
%		\end{tabular}
%		\caption{Four representations of $\Cyc_4$.}
%		\label{tbl:cyc4}
%	\end{table}
\end{example}
	
\begin{example}[$\Cyc_5$]
	The five fifth roots of unity are $\exp{\frac{2\pi im}{5}}$, $0 \leq m \leq 4$, so five representations of $\Cyc_5$ are presented in table \ref{table:Cyc5}. Note that $\rho_1^* = \rho_4$ and $\rho_2^* = \rho_3$.
	
%	\begin{table}[hbt!]
%		\begin{tabular}{c | c c c c c}
%			
%			$\Cyc_5$ & $e$ & $g$        & $g^2$      & $g^3$      & $g^4$      \\ \hline
%			$\rho_0$            & 1   & 1          & 1          & 1          & 1          \\
%			$\rho_1$            & 1   & $\omega^1$ & $\omega^2$ & $\omega^3$ & $\omega^4$ \\
%			$\rho_2$            & 1   & $\omega^2$ & $\omega^4$ & $\omega^1$ & $\omega^3$ \\
%			$\rho_3$            & 1   & $\omega^3$ & $\omega^1$ & $\omega^4$ & $\omega^2$ \\
%			$\rho_4$            & 1   & $\omega^4$ & $\omega^3$ & $\omega^2$ & $\omega^1$
%		\end{tabular}
%		\centering
%		\caption{Five representations of $\Cyc_5$. $\omega = e^{2 \pi i/5}$.}
%		\label{table:Cyc5}
%	\end{table}
\end{example}



\begin{table}[hbt!]
	\parbox[s]{.31\linewidth}{
		\centering
		\begin{tabular}{c | c c c}
			$\Cyc_3$ & $e$ & $g$        & $g^2$      \\ \hline
			$\rho_0$          & 1   & 1          & 1          \\
			$\rho_1$          & 1   & $\omega$   & $\omega^2$ \\
			$\rho_2$          & 1   & $\omega^2$ & $\omega$
		\end{tabular}
		\vfill
		\caption{Three reprs. of $\Cyc_3$. $\omega = e^{2 \pi i/3}$.}
		\label{table:Cyc3}
	}
	\hfill
	\parbox[s]{.31\linewidth}{
		\centering
		\begin{tabular}{c | c c cl}
			$\Cyc_4$ & $e$ & $g$  & $g^2$ & $g^3$ \\ \hline
			$\rho_0$           & 1   & 1    & 1     & 1     \\
			$\rho_1$           & 1   & $i$  & $-1$  & $-i$  \\
			$\rho_2$           & 1   & $-1$ & $1$  & $-1$   \\
			$\rho_3$           & 1   & $-i$ & $-1$   & $i$
		\end{tabular}
		\caption{Four reprs. of $\Cyc_4$.}
		\label{tbl:cyc4}
	}
	\hfill
	\parbox[s]{.31\linewidth}{
		\centering
		\begin{tabular}{c | c c c c c}
			$\Cyc_5$ & $e$ & $g$        & $g^2$      & $g^3$      & $g^4$      \\ \hline
			$\rho_0$            & 1   & 1          & 1          & 1          & 1          \\
			$\rho_1$            & 1   & $\omega$ & $\omega^2$ & $\omega^3$ & $\omega^4$ \\
			$\rho_2$            & 1   & $\omega^2$ & $\omega^4$ & $\omega$ & $\omega^3$ \\
			$\rho_3$            & 1   & $\omega^3$ & $\omega$ & $\omega^4$ & $\omega^2$ \\
			$\rho_4$            & 1   & $\omega^4$ & $\omega^3$ & $\omega^2$ & $\omega$
		\end{tabular}
		\caption{Five reprs. of $\Cyc_5$. $\omega = e^{2 \pi i/5}$.}
		\label{table:Cyc5}
	}
\end{table}



\subsubsection{Permutation representation}

	Given a group $G$, we chose a set $X$ that $G$ acts on by permutation. Let $V$ be a vector space spanned by a natural basis $(\bas_x)_{x \in X}$, then we have a representation
	\begin{align*}
		\rho:G \rightarrow \GL(V)
	\end{align*} 
	defined by its action on the basis vectors by, for any $g \in G$, 
	\begin{align*}
		\rho_g: \bas_x \mapsto \bas_{gx},
	\end{align*}
	that is $\rho$ inherited the group action of $G$ on $X$. It is a homomorphism since for any $g,h \in G$ we have
	\begin{align*}
		g \cdot (h \cdot \bas_x) &= g \cdot \bas_{hx} \\
		&= \bas_{ghx} \\
		&= (gh) \cdot \bas_x
	\end{align*}
	for any $x \in X$.%\marginnote{We have imposed a $G$-module structure on $V$.}
	In the $(\bas_x)_{x \in X}$-basis, $\rho_g$ are permutation matrices, which have one once in every row and column, and the rest of the entries are zero.
	
	Letting $G = \Sym_n$, it would be appropriate to choose the set $X = \{1, 2, \dots, n\}$ and to let any $\sigma \in \Sym_n$ permute any $1 \leq i \leq n$ by $i \mapsto \sigma(i)$. Choosing a basis $(\bas_i)_{i=1}^n$ to span a vector space $V$, we define a homomorphism
	\begin{align*}
		\rho: \Sym_n \rightarrow \GL(V)
	\end{align*}
	defined for any $\sigma \in \Sym_n$ as
	\begin{align*}
		\rho_\sigma: \bas_i \mapsto \bas_{\sigma(i)}
	\end{align*}
	for any $1 \leq i \leq n$. In this basis, the corresponding set of matrix representations are the permutation matrices
	\begin{align*}
		\rho_\sigma = (r_{ij})_{n \times n} \text{, where } r_{ij} = \delta_{j,\sigma(i)} = \begin{cases}
			1 \text{, if $j = \sigma(i)$,} \\
			0 \text{, otherwise,}
		\end{cases}
	\end{align*}
%	where $\delta_{k,l}$ is the Kronecker delta which vanishes if $k \neq l$ and is equal to one if $k=l$.
	in which the $i$th column has a 1 in the $\sigma(i)$th row, and the rest of the rows have a 0.
	
	Under this action, a vector
	\begin{align*}
		(a_1, a_2, \dots, a_n) = a_1 \bas_1 + a_2 \bas_2 + \dots + a_n \bas_n \in V
	\end{align*}
	is permuted to 
	\begin{align*}
		(a_{\sigma^{-1}(1)}, a_{\sigma^{-1}(2)}, \dots, a_{\sigma^{-1}(n)}) = a_1 \bas_{\sigma(1)} + a_2 \bas_{\sigma(2)} + \dots + a_n \bas_{\sigma(n)}  \in V.
	\end{align*}
	
%	\marginnote{Yet another interpretation is that $\sigma$ permutes the coordinate axes of a $n$-dimensional space.}
	
	The dimension of $V$ is $|X|$, for example the permutation representation of $\Sym_n$ is of degree $n$.
	
	\begin{example}[Permutation representation of $\Sym_2$]
		The symmetric group of degree 2 has two elements, $\Sym_2 = \{(1), (1,2)\}$ and their matrix representations in $\bas_1,\bas_2$-space are presented in Table~\ref{table:permS2}.
		\begin{table}[hbt!]
			\centering
			\begin{tabular}{c c}
				$\rho_{(1)} =
				\begin{pmatrix}
					1 & 0 \\ 0 & 1
				\end{pmatrix}$, &
				$\rho_{(1,2)} =
				\begin{pmatrix}
					0 & 1 \\ 1 & 0
				\end{pmatrix}$
			\end{tabular}
			\caption{Matrix representations of $\Sym_2$}
			\label{table:permS2}
		\end{table}
	\end{example}
	
	\begin{example}[Permutation representation of $\Sym_3$]\label{ex:permS3}
		Likewise, representations of $\Sym_3$ are presented in Table~\ref{table:permS3}.
		\begin{table}[hbt!]
			\centering
			\begin{tabular}{r r r}
				$\rho_{(1)} = 
				\begin{pmatrix}
					1 & 0 & 0 \\
					0 & 1 & 0 \\
					0 & 0 & 1
				\end{pmatrix}$, & 
				$\rho_{(1,2,3)} = 
				\begin{pmatrix}
					0 & 0 & 1 \\
					1 & 0 & 0 \\
					0 & 1 & 0
				\end{pmatrix}$, & 
				$\rho_{(1,3,2)} = 
				\begin{pmatrix}
					0 & 1 & 0 \\
					0 & 0 & 1 \\
					1 & 0 & 0
				\end{pmatrix}$, \\ & & \\
				$\rho_{(1,2)} = 
				\begin{pmatrix}
					0 & 1 & 0 \\
					1 & 0 & 0 \\
					0 & 0 & 1
				\end{pmatrix}$, &
				$\rho_{(1,3)} = 
				\begin{pmatrix}
					0 & 0 & 1 \\
					0 & 1 & 0 \\
					1 & 0 & 0
				\end{pmatrix}$, &
				$\rho_{(2,3)} = 
				\begin{pmatrix}
					1 & 0 & 0 \\
					0 & 0 & 1 \\
					0 & 1 & 0
				\end{pmatrix}$.
			\end{tabular}
			\caption{Matrix representations of $\Sym_3$}
			\label{table:permS3}
		\end{table}
	\end{example}
	
	\begin{example}[Permutation representation of $\Sym_4$]
		Yet again, matrix representations of some elements of $\Sym_4$ are presented in Table~\ref{table:permS4}. The elements chosen are one representative from every conjugacy class of $\Sym_4$.
		\begin{table}[hbt!]
			\centering
			\begin{tabular}{r r r}
				$\rho_{(1)} = \left(\begin{smallmatrix}
					1 & 0 & 0 & 0 \\
					0 & 1 & 0 & 0 \\
					0 & 0 & 1 & 0 \\
					0 & 0 & 0 & 1
				\end{smallmatrix}\right)$,  &
				$\rho_{(1,2)} = \left(\begin{smallmatrix}
					0 & 1 & 0 & 0 \\
					1 & 0 & 0 & 0 \\
					0 & 0 & 1 & 0 \\
					0 & 0 & 0 & 1
				\end{smallmatrix}\right)$, &
				$\rho_{(1,2)(3,4)} = \left(\begin{smallmatrix}
					0 & 1 & 0 & 0 \\
					1 & 0 & 0 & 0 \\
					0 & 0 & 0 & 1 \\
					0 & 0 & 1 & 0
				\end{smallmatrix}\right)$, \\ & & \\ &
				$\rho_{(1,2,3)} = \left(\begin{smallmatrix}
					0 & 0 & 1 & 0 \\
					1 & 0 & 0 & 0 \\
					0 & 1 & 0 & 0 \\
					0 & 0 & 0 & 1
				\end{smallmatrix}\right)$, &
				$\rho_{(1,2,3,4)} = \left(\begin{smallmatrix}
					0 & 0 & 0 & 1 \\
					1 & 0 & 0 & 0 \\
					0 & 1 & 0 & 0 \\
					0 & 0 & 1 & 0
				\end{smallmatrix}\right)$.
			\end{tabular}
			\caption{Some matrix representations of $\Sym_4$}
			\label{table:permS4}
		\end{table}
	\end{example}

\subsubsection{Regular representation}

	Following the same reasoning as in the previous section, but instead we let $G$ act on itself, ie. $X = G$. The corresponding vector space $V$ is spanned by the basis $(\bas_g)_{g \in G}$ constructed from the elements of $G$. A representation of $G$ in $V$ is then a map
	\begin{align*}
		\rho: G \rightarrow \GL(V)
	\end{align*}
	defined by
	\begin{align*}
		g \cdot \bas_h = \bas_{gh}
	\end{align*}
	for any $g,h \in G$.%\marginnote{We have created a Group Algebra $\CC[G]$.}
	
	The dimension of $V$ is $|G|$, for example the regular representation of $\Cyc_n$ is of degree $n$ and for $\Sym_n$ it is of degree $n!$, a number which grows increasingly quick for larger $n$, however the regular representation will be shown to be important later.
	
	\begin{example}[Regular representation of $\Cyc_3$]\label{example:regCyc3}
		For $G = \Cyc_3$ we calculate the action of $\Cyc_3$ on the $\bas, \ghat, \ghat_2$-basis thusly:
		\begin{align*}
			\begin{cases}
				e \cdot \bas = \bas, \\
				e \cdot \ghat = \ghat, \\
				e \cdot \ghat^2 = \ghat^2,
			\end{cases} \quad \begin{cases}
				g \cdot \bas = \ghat, \\
				g \cdot \ghat = \ghat^2, \\
				g \cdot \ghat^2 = \bas,
			\end{cases} \quad \begin{cases}
			g^2 \cdot \bas = \ghat^2, \\
			g^2 \cdot \ghat = \bas, \\
			g^2 \cdot \ghat^2 = \ghat,
			\end{cases}
		\end{align*}
		hence the permutation matrices in the $\bas, \ghat, \ghat_2$-basis are
		\begin{align*}
			\rho_e = \id, \quad \rho_g = \begin{pmatrix}
				0&0&1 \\ 1&0&0 \\ 0&1&0
			\end{pmatrix}, \quad \rho_{g^2} = \begin{pmatrix}
			0&1&0 \\ 0&0&1 \\ 1&0&0 
			\end{pmatrix}.
		\end{align*}
		Let's choose a new basis inspired by the third roots of unity 1, $\omega$ and $\omega^2$:
		\begin{align*}
			\begin{cases}
				\fbas_1 = \bas + \ghat + \ghat^2, \\
				\fbas_2 = \bas + \omega^2\ghat + \omega\ghat^2, \\
				\fbas_3 = \bas + \omega\ghat + \omega^2\ghat^2,
			\end{cases}
		\end{align*}
		where $\omega = e^{2 \pi i/3}$, corresponding to the change-of-basis matrix 
		\begin{align*}
			P = \begin{pmatrix}
				1&1&1 \\
				1&\omega^2&\omega \\
				1&\omega&\omega^2
			\end{pmatrix}, \quad \text{with inverse} \quad P^{-1} = \frac{1}{3}\begin{pmatrix}
				1&1&1 \\
				1&\omega&\omega^2 \\
				1&\omega^2&\omega
			\end{pmatrix}.
		\end{align*}
		Now, the matrices of $\Cyc_3$ in this new basis are
		\begin{align*}
			\rho_e &= P^{-1} \cdot \id \cdot P = \id, \\ 
			\rho_g &= P^{-1}\cdot\begin{pmatrix}
				0&0&1 \\ 1&0&0 \\ 0&1&0
			\end{pmatrix}\cdot P = \begin{pmatrix}
			1&0&0 \\ 0&\omega&0 \\ 0&0&\omega^2
			\end{pmatrix}, \\
			\rho_{g^2} &= P^{-1}\cdot\begin{pmatrix}
				0&1&0 \\ 0&0&1 \\ 1&0&0 
			\end{pmatrix}\cdot P = \begin{pmatrix}
			1&0&0 \\ 0&\omega^2&0 \\ 0&0&\omega
			\end{pmatrix}.
		\end{align*}
		We can clearly see that the regular representation of $\Cyc_3$ is the direct sum of the three degree 1 representations found in Table~\ref{table:Cyc3}, ie. $\rho_g = (1) \oplus (\omega) \oplus (\omega^2)$ after the change of basis.
	\end{example}
	
	\begin{example}[Regular representation of $\Sym_3$]\label{example:regSym3}
		Defining a homomorphism \begin{align*}
			\rho: S_3 \rightarrow \GL(V)
		\end{align*} on the vector space spanned by the basis\footnote{Explicitly, this is $\left\lbrace \bas_{(1)}, \bas_{(1,2)}, \bas_{(1,3)}, \bas_{(2,3)}, \bas_{(1,2,3)}, \bas_{(1,3,2)} \right\rbrace$.} $(\bas_{\sigma})_{\sigma \in \Sym_3}$ defined for every $\sigma$ by \begin{align*}
		\rho_\sigma \cdot \bas_\tau = \bas_{\sigma \circ \tau} = \bas_{\sigma\tau}
		\end{align*} for some $\tau \in \Sym_3$, we arrive at a degree 6 representation of $\Sym_3$. %It is clearly a larger degree representation of $S_n$ than any we have found so far.% Later we will find that for $\Sym_3$, $R$ is the direct sum of the trivial representation, the alternating representation and two copies of the standard representations, that is $S = T \oplus A \oplus 2S$.
		
		\textit{To be added: Table of $6 \times 6$-matrices, one per conjugacy class for example. Also: Show that perm.repr. of $\Sym_3$ is completely decomposed to trivial + alternating + two copies of standard, especially since the multiplicity of the standard representation is $>1$.}
	\end{example}
	
	

%\subsubsection{Coset representations}
%
%Let $H$ be a subgroup of $G$. Let $\mathcal{H} = \{g_1H, g_2H, \dots, g_kH\}$ be the set of left cosets of $H$ in $G$ and define the action of $G$ on $\mathcal{H}$ by $g(g_iH) = (gg_i)H$ for any $g\in G$ and $g_iH \in \mathcal{H}$. The $g_i$ are called transversals of their respective $g_iH$. Construct a vector space $V$ by spanning it with the transversals as a basis, denoted $(\ghat_i)_{i=1}^k$ and define the group action of $G$ on $V$ by replicating the group action of $G$ on $\mathcal{H}$.\cite[Example 1.3.5.]{Sagan}
%
%\begin{example}[Trivial representation again]
%	Choosing $H = G$ (the non-proper subgroup), the coset representation reduces to the trivial representation.
%\end{example}
%
%\begin{example}[Regular representation again]
%	Choosing $H = \{e\}$ (the trivial subgroup), then $\mathcal{H} = G$ and we construct a vector space $V$ using the elements of $G$ as a basis, hence we have found the regular representation.
%\end{example}
%
%\begin{example}[Choosing a smarter subgroup of $\Sym_3$]
%	As laid out in~\cite[Example 1.3.5.]{Sagan}, choosing the subgroup of $\Sym_3$ to be $H = \{(1), (2,3)\}$, then the cosets of $H$ are $\mathcal{H} = \{ H, (1,2)H, (1,3)H \}$, and computing the matrices will yield the permutation representation of $\Sym_3$.
%\end{example}
\subsection{Subrepresentations}\label{sect:subrepr}

As laid out in \cite[Sect.?.?.]{FultonHarris}, we are looking for representations which are said to be ``atomic'' and conversely, for any arbitrary representation we wish to find how it is composed of these indecomposable representations. To proceed we need the notion of a vector space map that respects or conserves the group action.	

\begin{definition}[$G$-linear map]\cite{FultonHarris}
	Let $V$ and $W$ be two representation spaces of a group $G$. A vector space map 
	\begin{align*}
		\varphi: V \rightarrow W
	\end{align*} 
	is called a $G$-linear map if it commutes with the group action of $G$, ie. for any $\vvec \in V$ and $g \in G$ we have
	\begin{align*}
		\varphi( g \cdot \vvec ) = g \cdot \varphi (\vvec),
	\end{align*}
	or in terms of the maps $\rho_g^V: G \rightarrow \GL(V)$ and $\rho_g^W: G \rightarrow \GL(W)$,
	\begin{align*}
		\varphi \circ \rho^V_g (\vvec) = \rho^W_g \circ \varphi (\vvec).
	\end{align*}
	Equivalently one can say that the diagram in Figure~\ref*{fig:commapmorp} is commutes for every $g \in G$.
	\begin{figure}[hbt!]
		\centering
%		\[\begin{tikzcd}
%			V \arrow{r}{\varphi} \arrow[swap]{d}{g \cdot} & W \arrow{d}{g \cdot} \\
%			V \arrow{r}{\varphi} & W
%		\end{tikzcd}\]
		% https://q.uiver.app/#q=WzAsNCxbMCwwLCJWIl0sWzIsMCwiVyJdLFswLDIsIlYiXSxbMiwyLCJXIl0sWzAsMiwiZ1xcY2RvdCIsMl0sWzEsMywiZ1xcY2RvdCJdLFsyLDMsIlxcdmFycGhpIiwyXSxbMCwxLCJcXHZhcnBoaSJdXQ==
		\[\begin{tikzcd}
			V && W \\
			\\
			V && W
			\arrow["\varphi", from=1-1, to=1-3]
			\arrow["g\cdot"', from=1-1, to=3-1]
			\arrow["g\cdot", from=1-3, to=3-3]
			\arrow["\varphi"', from=3-1, to=3-3]
		\end{tikzcd}\]
		\caption{The map $\varphi$ is $G$-linear if the diagram commutes.}
		\label{fig:commapmorp}
	\end{figure}
\end{definition}

Consider the case where $W$ is a $G$-invariant\footnote{For $W$ to be a $G$-invariant space, it means that for every $g \in G$, $g \cdot \wvec \in W$ for any $\wvec \in W$ or that $\rho_g(W) \subseteq W$ for any $g\in G$.} vector subspace of $V$% and let $\pi:V \rightarrow W$ be the projection map of $V$ onto $W$. 
and let $\rho$ be a representation of $G$ in $V$. Then the restriction of $\rho$ to $W$, here denoted $\rho|_W$, is a isomorphism of $W$ onto itself since $W$ is $G$-invariant,
%(the action of $g$ on any $\wvec \in W$ will return the same $\wvec$), 
and thus $\rho_{gh}|_W = \rho_g|_W \cdot \rho_h|_W$, hence $\rho|_W: G \rightarrow \GL(W)$ is a representation of $G$ in the subspace $W$ of $V$, motivating the following definition:

\begin{definition}[Subrepresentation]\label{def:subrepr}
	Let $G$ be a finite group and $\rho$ be a representation of $G$ in a vector space $V$. A restriction of $\rho$ to a $G$-invariant vector subspace $W$ of $V$ is called a subrepresentation of $\rho$.
\end{definition}

%In other words, a subrepresentation is a $G$-invariant vector subspace of a ``parent'' representation. 

\begin{example}[Trivial subspaces]
	Any representation has itself as well as the zero space as subrepresentations.
\end{example}

\subsubsection{Some proper and non-trivial subrepresentations}

Let $G$ be any group and let us study its permutation and regular representations. 

\begin{example}[Trivial representation inside the permutation representation]\cite[Example 1.4.3.]{Sagan}\label{example:trivpermrepr}
	Let $G$ act on a set $X = \{x_1, x_2, \dots x_k\}$, where $k = |X|$, and let $V$ be the vector space spanned by the basis $(\xhat_1, \xhat_2, \dots, \xhat_k)$. Consider the one-dimensional subspace of $V$ spanned by the sum of all basis vectors, ie. 
	\begin{align*}	
		W = \Span\{ \xhat_1 + \xhat_2 + \dots + \xhat_k \}.
	\end{align*}A vector $\wvec \in W$ is a scalar multiple of this sum, and for any $g \in G$, the action of $g$ on $\wvec$ will simply reorder this sum and return the same $\wvec$, that is for any $g \in G$, we have that $\rho_g|_W = 1$. Thus $W$ is a $G$-invariant subspace of $V$ and the permutation representation has the trivial representation as a subrepresentation.
\end{example}

\begin{example}[Trivial representation inside the regular representation]\cite[Example 1.4.4.]{Sagan}
	Similarly to the last example, we span a vector space $V$ by a basis $(\ghat_1, \ghat_2, \dots, \ghat_{k})$ for every $g_i \in G$, where $k = |G|$, and consider the subspace 
	\begin{align*}
		W = \Span\{ \ghat_1 + \ghat_2 + \dots + \ghat_k \}. 
	\end{align*}
	Completely analogously to the last example, $W$ is shown to be a $G$-invariant subspace of $V$ and the regular representation also has the trivial representation as a subrepresentation.
\end{example}

\begin{example}[Alternating representation inside the regular representation of $\Sym_n$]\cite[Example 1.4.4.]{Sagan}
	Let $G = \Sym_n$, then for the regular representation, $V$ is spanned by a basis vector for every $\sigma \in \Sym_n$. Let $W$ be the subspace of $V$ spanned by the sum $\sum_{\sigma \in \Sym_n} \sgn(\sigma) \bas_{\sigma}$. The action of a $\tau \in \Sym_n$ on a $\wvec \in W$ is
	\begin{align*}
		\tau (\wvec) &= \tau \cdot a\sum_{\sigma \in \Sym_n} \sgn(\sigma) \bas_{\sigma} \\
		&= a \sum_{\sigma \in \Sym_n} \sgn(\sigma) \bas_{\tau\sigma} \\
		&= a \frac{1}{\sgn\tau} \sum_{\sigma \in \Sym_n} \sgn(\tau\sigma) \bas_{\tau\sigma} 
	\end{align*}
	where $a \in \CC$. Let $\pi = \tau\sigma$, then the sum can instead be over all these $\pi \in \Sym_n$, hence $W$ is a $G$-invariant subspace of $V$ and the alternating representation is also found inside the regular representation of the symmetric group.
\end{example}

\subsubsection{Subrepresentations as a kernel}

From linear algebra, we are familiar with the kernel and image of a map (see Definition~\ref{def:kernelimage}).

\begin{proposition}[Kernel and image of a $G$-linear map]\label{prop:kernelimagelinearmap}
	Let $V$ and $W$ be representation spaces of a group $G$ and let $\varphi: V \rightarrow W$ be a $G$-linear map. Then the kernel and the image of $\varphi$ are also representations of $G$ and more specifically, $\ker \varphi$ is a subrepresentation of $V$ and $\im \varphi$ is a subrepresentation of $W$.
\end{proposition}
\begin{proof}
	\begin{itemize}
		\item[i)] Let $\vvec \in \ker \varphi$. Then $g \cdot \varphi (\vvec) = \0$ since by definition $\varphi(\vvec) = \0$, but then we must also have $\varphi ( g \cdot \vvec ) = \0$ since $\varphi$ is a $G$-linear map, implying that $g \cdot \vvec \in \ker \varphi$. Since this holds for any $g \in G$ and $\vvec \in \ker \varphi$, the kernel of $\varphi$ is a $G$-invariant subspace of $V$ and thus a subrepresentation of $V$. %\marginnote{Subrepr. not yet defined}
		\item[ii)] For a $\wvec \in \im \varphi$, by definition there exists a $\vvec \in V$ such that $\varphi(\vvec) = \wvec$. Since $g \cdot \vvec \in V$, then we have that $\varphi ( g \cdot \vvec) \in W$. This implies that $g \cdot \varphi (\vvec)$ is also in $W$. Since this holds for any $g \in G$ and $\vvec \in V$, the image of $\varphi$ is a $G$-invariant subspace of $W$ and a thus a subrepresentation of $W$.	\qedhere
		%			\item[i)] Let $\vvec \in \ker \varphi$. Then $\rho_g \circ \varphi (\vvec) = \0$, but then we must also have $\varphi \circ \rho_g (\vvec) = \0$, implying that $\rho_g (\vvec) \in \ker \varphi$. Since this holds for any $g \in G$ and any $\vvec \in \ker \varphi$, the kernel of $\varphi$ is a $G$-invariant subspace of $V$ and thus a subrepresentation of $V$.
		%			\item[ii)] For a $\wvec \in \im \varphi$, there exists a $\vvec \in V$ such that $\varphi(\vvec) = \wvec$. Since $\rho_g(\vvec) \in V$, then we have that $\varphi \circ \rho_g(\vvec) \in W$. This implies that $\rho_g \circ \varphi (\vvec)$ is also in $W$. Since this holds for any $g \in G$ and $\vvec \in V$, the image of $\varphi$ is a $G$-invariant subspace of $W$ and a thus a subrepresentation of $W$.	\qedhere
	\end{itemize}
\end{proof}

Recall from linear algebra the existence of complementary vector subspaces of a vector space. Is there a similar property of subrepresentations? Consider the case when $W$ is a subrepresentation of $V$ and let $\pi: V \rightarrow W$ be the projection of $V$ onto $W$. Recall (Corollary~\ref{thm:compsubspaces})~that for each subspace $W\leq V$ there exists a complementary and disjoint subspace $W' \leq V$ such that $V = W \oplus W'$. This means that every $\vvec \in V$ is partitioned into $\wvec \oplus \wvec'$ where $\wvec \in W$ and $\wvec' \in W'$ such that 
\begin{align*}
	\begin{cases}
		\pi(\vvec) = \wvec, \\
		\pi(\wvec) = \wvec \text{ and } \\
		\pi(\wvec') = \0.
	\end{cases}
\end{align*} 
The image of the projection is clearly $W$ and the kernel is $W'$. But is $\pi$ a $G$-linear map?
\begin{proposition}[Existence of complementary subrepresentations]\label{prop:complsubrepr}
	Let $\rho$ be a representation of a finite group $G$ in a finite-dimensional vector space $V$. Let $W$ be a $G$-invariant subspace of $V$. Then there exists a $G$-invariant subspace $W'$ of $V$ complementary to $W$ such that $V = W \oplus W'$.
\end{proposition}
\begin{proof}
	\textit{THIS PROOF MISSES SOME PARTS: WE NEED TO SHOW THAT $\barpi$ IS STILL A PROJECTION $V \rightarrow W$.}
	By Proposition~\ref{prop:kernelimagelinearmap} we know that a $G$-linear map from one representation space $V$ to another $W$ has its kernel and image as subrepresentations of $V$ and $W$ respectively. If we can find such a $G$-linear map from $V$ to a $G$-invariant subspace $W \leq V$ we are done.
	The projection $\pi$ may not generally be a $G$-linear map.% since $\pi(g \cdot \vvec)$ is not generally the same as $g \cdot \pi(\vvec)$. 
	\textit{NEEDS A COUNTEREXAMPLE?}
%	for a $\vvec \in V$ we have that \marginnote{There is a mistake here}
%	\begin{align*}
%		\pi(g \cdot \vvec) = \pi(\wvec) = \wvec,
%	\end{align*}
%	but 
%	\begin{align*}
%		g \cdot \pi(\vvec) = g \cdot \wvec,
%	\end{align*}
%	and the equality $g \cdot \wvec = \wvec$ is not generally true for every $g \in G$.
	Instead, consider taking the \textit{average} of $\pi$ over $G$,
	\begin{align*}
		\barpi = \frac{1}{|G|} \sum_{g \in G} \rho_g \cdot \pi \cdot \rho_g^{-1},
	\end{align*}
	and see if it conserves the action of $G$. That is, does $\barpi\cdot \rho_g = \rho_g \cdot \barpi$ hold for every $g \in G$? Equivalently, let's consider
	\begin{align*}
		\rho_g \cdot \barpi \cdot \rho_g^{-1} &= \frac{1}{|G|} \sum_{h \in G} \rho_g \cdot \rho_h \cdot \pi \cdot \rho_h^{-1} \cdot \rho_g^{-1} & \text{(Def. of $\barpi$)}\\
		&= \frac{1}{|G|} \sum_{h \in G} \rho_{gh} \cdot \pi \cdot \rho_{gh}^{-1} & \text{($\rho$ homom., $(gh)^{-1}=h^{-1}g^{-1}$)}\\
		&= \frac{1}{|G|} \sum_{g' \in G} \rho_{g'} \cdot \pi \cdot \rho_{g'}^{-1} & \text{(Let $g' = gh$)} \\
		&= \barpi.
	\end{align*}
	So, $\barpi \cdot \rho_g = \rho_g \cdot \barpi$ for every $g \in G$ and thus $\barpi: V \rightarrow W$ is a $G$-linear map, which means that its kernel, $\ker \barpi$, identified to be the complement $W'$ of $W$ in $V$, is also a $G$-invariant subspace of $V$ and thereby a subrepresentation of $V$.
\end{proof}

\subsubsection{The standard represention of $\Sym_n$}

In Example~\ref{example:trivpermrepr}, the trivial representation was found to be a subrepresentation of the permutation representation. Let $G= \Sym_n$ and let $V$ be the permutation representation and assign it the basis $(\bas_i)_{i=1}^n$. The one-dimensional subspace $W$ spanned by the sum of all $\bas_i$ was found to be the trivial representation, then by Proposition~\ref{prop:complsubrepr}, there exists another subrepresentation $W'$ of $V$ complementary to $W$. The dimension of $W'$ is $n-1$. Let's introduce to $V$ a inner product $(\cdot|\cdot)$ that fulfills the expected properties, in particular:
\begin{align*}
	(\bas_i|\bas_j) = \delta_{i,j} = \begin{cases}
		1 \text{, if $i=j$,}\\
		0 \text{, if $i \neq j$.}
	\end{cases}
\end{align*}
for the basis $(\bas_i)_{i=1}^n$. Then the basis $(\fbas_j)_{j=1}^{n-1}$ of $W'$ can be constructed by ensuring that each $\fbas_j$ is orthogonal to the basis vector of $W$, that is
\begin{align*}
	\left( \fbas_j \ \middle\vert \ \sum_{i=1}^n \bas_i \right) = 0
\end{align*}
for every $1 \leq j \leq n-1$. One can see that one such basis is found by for example choosing every basis vector $\fbas_j$ to be the difference of two sequent $\bas_i$, that is choosing the basis $(\bas_j - \bas_{j+1})_{j=1}^{n-1}$ for $W'$. This is called the \emph{Standard representation of $\Sym_n$.}

\begin{example}[Standard representation of $\Sym_3$]
	The matrices of the standard representation is found by studying the action inherited from $\Sym_3$ on the basis 
	\begin{align*}
		\begin{cases}
			\fbas_1 = \bas_1-\bas_2, \\
			\fbas_2 = \bas_2-\bas_3.
		\end{cases}
	\end{align*}
	For example, $(1,2)\cdot (\bas_1-\bas_2) = \bas_2-\bas_1 = -\fbas_1$ and $(1,2) \cdot (\bas_2-\bas_3) = \bas_1-\bas_3 = \fbas_1 + \fbas_2$, hence $\rho_{(1,2)} = \left(\begin{smallmatrix}
		-1 & 1 \\ 0 & 1
	\end{smallmatrix}\right)$. The matrices are presented in Table~\ref{table:standardS3}.
	
	\begin{table}[hbt!]
		\centering
		\begin{tabular}{r r r}
			$\rho_{(1)} = 
			\begin{pmatrix}
				1 & 0 \\
				0 & 1
			\end{pmatrix}$, & 
			$\rho_{(1,2,3)} = 
			\begin{pmatrix}
				0 & -1 &\\
				1 & -1
			\end{pmatrix}$, & 
			$\rho_{(1,3,2)} = 
			\begin{pmatrix}
				-1 & 1  \\
				-1 & 0
			\end{pmatrix}$, \\ & & \\
			$\rho_{(1,2)} = 
			\begin{pmatrix}
				-1 & 1 \\
				0  & 1
			\end{pmatrix}$, &
			$\rho_{(1,3)} = 
			\begin{pmatrix}
				-1 & 0 \\
				0 & -1 
			\end{pmatrix}$, &
			$\rho_{(2,3)} = 
			\begin{pmatrix}
				1 & 0 \\
				1 & -1 
			\end{pmatrix}$.
		\end{tabular}
		\caption{Matrices of the standard representations of $\Sym_3$}
		\label{table:standardS3}
	\end{table}
\end{example}

\subsection{Irreducible representations}

In Section~\ref{sect:subrepr} we found that a given representation can be divided into complementary subrepresentations, now we introduce the notion of an ``indivisible'' representation.

\begin{definition}[Irreducible representation]
	If there are no proper and non-trivial $G$-invariant subspaces of a representation $V$, then it is said to be \textit{irreducible}. 
\end{definition}

We have already met a few of these.

\begin{example}[Degree 1 representations are irreducible]\cite[Example 1.4.2.]{Sagan}
	A vector space of dimension 1 has no other subspace other than itself and the zero space, thus it is irreducible. Hence the trivial representation of degree 1 of any group and the alternating representation of $\Sym_n$ discussed in Section~\ref{sect:basicreprs} and the degree 1 representations of $\Cyc_n$ found in Section~\ref{sect:deg1cycn} are irreducible representations of their respective groups.
\end{example}

The results from Proposition~\ref{prop:complsubrepr} invites the notion of \textit{complete reducibility} of an arbitrary representation. This is presented in the following theorem.

\begin{theorem}[Maschke's theorem]\label{thm:maschkes}
	Let $G$ be a finite group and let $V$ be any representation space of $G$. Then $V$ is composed as a direct sum of a number of not necessarily distinct subrepresentations $W_i$ of $V$, or in other words,
	\begin{align*}
		V = \bigoplus_i W_i
	\end{align*}
	\begin{note}
		The direct sum can be expressed in terms of distinct subrepresentation $W_i$ of $V$ and their multiplicities $a_i$ as
		\begin{align*}
			V = \bigoplus_i a_i W_i
		\end{align*}
	\end{note}
\end{theorem}

\begin{example}[Decomposition of permutation representation of $\Sym_n$]
	Let $G = \Sym_n$, and let $P$ be the permutation representation space, $S$ be the standard representation space and $T$ be the trivial representation space, then $V$ is fully decomposed as 
	\begin{align*}
		P = T \oplus V.
	\end{align*}
\end{example}

%\begin{example}[Decomposition of regular representation of $\Sym_n$]
%	Again, let $G = \Sym_n$, and let $R$ be the regular representation space, 
%\end{example}

\begin{theorem}[Schur's lemma]
	content...
\end{theorem}

%	In Example~\ref{example:comptrivpermrepr}, a subrepresentation of the permutation representation complementary to the trivial representation was found. Show it is irreducible? 

%	Consider the action of a group $G$ on the basis $(\bas_i - \bas_{i+1})_{i=1}^n$. For a $g \in G$ we have
%	\begin{align*}
	%		g \cdot (\bas_i - \bas_{i+1}) &= \bas_{g(i)} - \bas_{g(i+1)}
	%	\end{align*}
%	
%	\begin{example}
	%		hello
	%	\end{example}
	


		
	

%The section with references.

\clearpage{\thispagestyle{empty}}
\bibliographystyle{alphaurl} %This style prints the urls of online sources.


% All references are to be stored in the file bibliographyFile.bib! 
% It is in the same folder as your main .tex file.
% It is using the bibtex syntax, and you can get these entries from google scholar, etc.
% Your supervisor should know,
\bibliography{bibFile}
%\bibliography{bibliographyFile} 

\end{document}
